\documentclass[11pt]{article}
\usepackage[margin=2.5cm]{geometry}
\usepackage{amsmath,amssymb,amsthm}
\usepackage{booktabs}
\usepackage{array}
\usepackage{hyperref}

\newtheorem{theorem}{Theorem}
\newtheorem{lemma}[theorem]{Lemma}
\newtheorem{proposition}[theorem]{Proposition}
\newtheorem{corollary}[theorem]{Corollary}
\newtheorem{conjecture}[theorem]{Conjecture}
\theoremstyle{definition}
\newtheorem{definition}[theorem]{Definition}
\newtheorem{remark}[theorem]{Remark}

\newcommand{\F}{\mathbb{F}}
\DeclareMathOperator{\Alt}{Alt}
\DeclareMathOperator{\rad}{rad}
\DeclareMathOperator{\Tr}{Tr}
\DeclareMathOperator{\rk}{rank}

\title{Quadratic Differentially 4-Uniform Permutations\\ in Dimensions Divisible by Four:\\
A Desarguesian Obstruction and Non-Existence in Dimension 8}
\author{Ali Huseynli\\ \small Azerbaijan Technical University, Baku, Azerbaijan}
\date{}

\begin{document}
\maketitle

\begin{abstract}
Quadratic S-boxes are attractive for masked implementations because they admit three-share
threshold implementations without decomposition into several stages, and differential
uniformity~$4$ is the best possible value
for a quadratic permutation of even dimension. By a result of Charpin and Peng, a quadratic
differentially 4-uniform permutation of $\F_2^n$ is differentially two-valued and all its
components have rank $n-2$. All known examples of such permutations live in dimensions
$n\equiv 2\pmod 4$. We study
the case $4\mid n$. We observe that the kernels of the derivatives of such a permutation form a
line spread of $\F_2^n$ which is totally isotropic for all components and is mapped onto a line
spread of the codomain. We show that this spread cannot be Desarguesian when
$4\mid n$; more generally, a quadratic permutation whose derivative kernels form a Desarguesian
$t$-spread exists only if $n/t$ is odd, which contains the classical fact that quadratic APN
permutations exist only in odd dimension. As a consequence no $\F_4$-quadratic map, for instance
no map $x^{2^k+1}+\ell(x)$ with $k$ even and $\ell$ linear, is a differentially 4-uniform permutation
when $4\mid n$ (for power maps alone this is classical). For
$n=8$ we settle the question: by an exhaustive computer search over the possible spaces of
components, cross-checked by a second, separately written implementation, there is no quadratic differentially
4-uniform permutation of $\F_2^8$; more generally every quadratic $(8,8)$-function whose
nonzero derivatives are all $4$-to-$1$ has a bent component. The search rests on a simple
closure property of the radicals of the components.
\end{abstract}

\noindent\textbf{Keywords:} differential uniformity, quadratic permutation, S-box, spread,
alternating bilinear form, constant-rank subspace, computer-assisted proof.

\noindent\textbf{MSC (2020):} 94A60, 94D10, 11T71, 15A63, 51E23.

\section{Introduction}
Masked implementations of block ciphers pay for every non-linear operation. A threshold
implementation of a function of algebraic degree $d$ needs at least $d+1$ shares
\cite{NRR06}; for a quadratic S-box three shares suffice for a correct and non-complete sharing
computed in one stage, whereas S-boxes of higher degree need more shares or a decomposition
into several stages of lower degree separated by registers \cite{BNNRS15}. (Uniformity of the
shared function is a separate requirement, which may need fresh randomness or a particular
sharing; a low algebraic degree alone does not make an implementation secure.) So the
cryptographic quality achievable by quadratic permutations is a natural design limit. For a permutation of $\F_2^n$ with $n$ even
the differential uniformity is at least $4$ when the permutation is quadratic
(Remark~\ref{rem:du-at-least-4}); the value $4$ is attained by the Gold permutations
$x^{2^k+1}$ with $\gcd(k,n)=2$ precisely when $n\equiv 2\pmod 4$ \cite{Gold68,Nyb93}, and
other known quadratic examples, such as the binomials of Bracken, Tan and Tan \cite{BTT12}
and the Dembowski--Ostrom quadrinomials over $\F_{2^{2m}}$ with $m$ odd studied in
\cite{KMCLLJ22}, also live in dimensions $n\equiv 2\pmod 4$. For $n\le6$ all quadratic
permutations are classified up to affine equivalence \cite{DMB19}; for $n=6$ eight of the
classes have differential uniformity $4$. For $4\mid n$ differentially
4-uniform permutations do exist, for instance the inverse map \cite{Nyb93} and the cubic power
map $x^{2^{2k}+2^k+1}$ for $n=4k$ with $k$ odd \cite{BL10}, but no quadratic one is known. We
show that none exists for $n=8$.

The same dichotomy between $n\equiv 2$ and $n\equiv 0\pmod 4$ appears in the decomposition
problem for permutations into quadratic ones \cite{NNR19}: if $4\nmid n$, every permutation of
$\F_{2^n}$ is a composition of quadratic power permutations and affine permutations $x\mapsto
ax+b$, whereas if $4\mid n$, no power permutation with differential uniformity below $16$ is a
composition of quadratic and linear power permutations \cite{ACPBSN25}. Quadratic
permutations with differential uniformity $4$ in dimension $4\mid n$ are therefore of interest
both as candidate S-boxes and as building blocks of such decompositions.

\paragraph{Question.} Does $\F_2^n$ with $4\mid n$ admit a quadratic permutation with
differential uniformity $4$? The first open case was $n=8$, the dimension of the AES S-box
(which is based on the inverse map and is not quadratic); we answer it negatively (Theorem~\ref{thm:n8}).

\paragraph{Known structure.} Charpin and Peng \cite[Theorems~1 and~7]{CP19} proved that a
quadratic differentially 4-uniform permutation of $\F_2^n$ ($n$ even) is differentially
two-valued $\{0,4\}$, has no bent component and has nonlinearity $2^{n-1}-2^{n/2}$ with a
single Walsh amplitude; in the language of alternating forms, every component has rank exactly
$n-2$ and every derivative kernel has dimension $2$ (Theorem~\ref{thm:constrank} below). Their
proof rests on the identity \cite[Theorem~5]{CP19} which we recall as
Lemma~\ref{lem:count}.

\paragraph{Contributions.}
\begin{enumerate}
\item Kernel spreads (Section~\ref{sec:structure}): the derivative kernels of a quadratic DU-4
permutation form a line spread that is totally isotropic for all components; the restriction
of the permutation to each kernel line is linear, and the images form a line spread of the
codomain. This gives a partial normal form. (Isotropy of each single kernel is \cite[Lemma~1]{MTX19}.)
\item The Desarguesian obstruction (Section~\ref{sec:f4}), our first main result: a quadratic
permutation all of whose derivative kernels have dimension $t$ and form a Desarguesian
$t$-spread exists only if $n/t$ is odd (Theorem~\ref{thm:general-t}). For $t=1$ this is the
classical fact that quadratic APN permutations exist only in odd dimension; for $t=2$ it shows
that, if $4\mid n$, the kernel spread of a quadratic DU-4 permutation is not Desarguesian. This
excludes all $\F_4$-quadratic maps, for example Gold maps $x^{2^k+1}$ with $k$ even plus
arbitrary linear maps (for the power maps themselves the conclusion is classical), and
covers the case $n=4$.
\item Non-existence for $n=8$ (Section~\ref{sec:n8search}), our second main result: there is no
quadratic DU-4 permutation of $\F_2^8$ (Theorem~\ref{thm:n8}). By the result of Charpin and
Peng it suffices to show that $\Alt(8,\F_2)$ has no $8$-dimensional subspace of constant rank
$6$ in which every point lies in exactly three radicals. We prove this by an exhaustive search
built on a radical-closure lemma (Lemma~\ref{lem:closure}), whose completeness is proved in
Appendix~\ref{app:search}; the search takes a few CPU-hours and was cross-checked by a second,
separately written implementation and a certificate for its symmetry reductions.
\item Local structure of kernel spreads (Section~\ref{sec:local}): for every even $n\ge6$, two kernel
lines span a $4$-space containing $2$, $3$ or $5$ kernel lines, and the dimension of
$\langle\beta(L,M)\rangle$ is determined by this number (Proposition~\ref{prop:local}). We also
summarise earlier partial results for $n=8$ obtained with other methods, among them that the
quadratic part of such a permutation could have no symmetry.
\item A necessary condition in terms of constant-rank subspaces of alternating forms, and a
stronger conjecture that would settle the question (Section~\ref{sec:constrank}).
\end{enumerate}

\section{Preliminaries}\label{sec:prelim}
Let $V=\F_2^n$. For $F:V\to V$ and $b\in V\setminus\{0\}$ we write
$F_b(x)=\langle b,F(x)\rangle$ for the components of $F$. $F$ is quadratic if its algebraic
degree is at most $2$, that is, if all its second-order derivatives are constant (affine maps
are included); we normalise $F(0)=0$ and write
\[
\beta(u,v)=F(u+v)+F(u)+F(v),
\]
a symmetric bi-additive map with $\beta(u,u)=0$. For each $b\neq 0$,
$B_b(u,v)=\langle b,\beta(u,v)\rangle$ is an alternating bilinear form with radical
$R_b=\rad B_b$; $\rk B_b$ is even. For $a\neq 0$ put
$K_a=\{v:\beta(a,v)=0\}\ni a$ and $s_a=\dim K_a\ge 1$. Since
$D_aF(x)=F(x+a)+F(x)=\beta(a,x)+F(a)$, the differential uniformity is
$\delta_F=\max_{a\neq 0}2^{s_a}$. A Boolean function is \emph{bent} if its Walsh transform
has constant absolute value $2^{n/2}$; a quadratic component is bent iff $\rk B_b=n$, and
semi-bent (Walsh values in $\{0,\pm 2^{(n+2)/2}\}$, $n$ even) iff $\rk B_b=n-2$. Bent functions
are never balanced \cite{Carlet21}, while every component of a permutation is balanced.
Composing $F$ with affine permutations on either side and adding affine maps (EA-equivalence)
preserves $\delta_F$ and the ranks of the components; composing with affine permutations
preserves bijectivity. When we identify $V$ with $\F_{2^n}$ (Section~\ref{sec:f4}), components
are indexed through the absolute trace pairing: $F_b(x)=\Tr_{\F_{2^n}/\F_2}(bF(x))$ and
$B_b(u,v)=\Tr_{\F_{2^n}/\F_2}(b\,\beta(u,v))$.

\section{Structure of quadratic differentially 4-uniform permutations}\label{sec:structure}

\begin{lemma}[Counting identity {\cite[Theorem~5]{CP19}}]\label{lem:count}
For every quadratic $F$,
\[\sum_{a\neq 0}\bigl(2^{s_a}-1\bigr)=\sum_{b\neq 0}\bigl(2^{\,n-\rk B_b}-1\bigr).\]
\end{lemma}
\begin{proof}
Both sides count the pairs $(a,b)$ of nonzero vectors with $B_b(a,\cdot)=0$: for fixed
$a$ these $b$ form $\{b:\langle b,\beta(a,v)\rangle=0\ \forall v\}=(\mathrm{Im}\,\beta(a,\cdot))^{\perp}$,
of dimension $n-(n-s_a)=s_a$; for fixed $b$ these $a$ form $R_b\setminus\{0\}$.
\end{proof}
\begin{remark}\label{rem:du-at-least-4}
If $n$ is even and $F$ is a quadratic permutation then no component is bent (bent
functions are not balanced), so $\rk B_b\le n-2$ and $|R_b|\ge 4$ for all $b$;
Lemma~\ref{lem:count} then gives $\sum_a (2^{s_a}-1)\ge 3(2^n-1)$, hence $\delta_F\ge 4$
(cf.\ \cite[Theorem~1]{CP19}).
\end{remark}

\begin{remark}\label{rem:count-general}
The same double count applies to the restriction of $\beta$ to $U\times U$ for any subspace $U$:
$\sum_{a\in U\setminus\{0\}}(2^{\,n-\dim\beta(a,U)}-1)=\sum_{b\neq0}(2^{\,\dim U-\rk(B_b|_U)}-1)$.
\end{remark}
\begin{remark}
In \cite{CP19} the identity is written as $\sum_{\lambda\neq0}2^{\ell(\lambda)}=\sum_{a\neq0}2^{d(a)}$,
where $\ell(\lambda)=\dim R_\lambda$ and $d(a)=s_a$; subtracting $2^n-1$ from both sides gives
the form above. For quadratic APN functions ($s_a=1$) it gives the classical fact that in even
dimension such functions have bent components \cite[Remark~3]{CP19}.
\end{remark}

The following theorem is due to Charpin and Peng \cite[Theorem~1, Lemma~4 and
Theorem~7]{CP19}. We include the short derivation from Lemma~\ref{lem:count} because the
equality case is used repeatedly below.
\begin{theorem}[\cite{CP19}]\label{thm:constrank}
Let $n$ be even and $F$ a quadratic permutation of $\F_2^n$ with $\delta_F=4$. Then
$\rk B_b=n-2$ for every $b\neq 0$ and $s_a=2$ for every $a\neq 0$.
\end{theorem}
\begin{proof}
By Remark~\ref{rem:du-at-least-4} the right-hand side of Lemma~\ref{lem:count} is at least
$3(2^n-1)$, with equality iff all ranks equal $n-2$; $\delta_F=4$ gives $s_a\le 2$, so the
left-hand side is at most $3(2^n-1)$, with equality iff all $s_a=2$.
\end{proof}
\begin{remark}
In the terminology of \cite{CP19}, $F$ is differentially two-valued $\{0,4\}$ and plateaued with
single amplitude $2^{(n+2)/2}$ \cite[Theorem~7]{CP19}. For a quadratic function without bent
components, Lemma~\ref{lem:count} shows that the first property implies the second (every term
on the right is at least $3$ and the left-hand side is $3(2^n-1)$), and that the second implies
the first if $\delta_F\le4$. Whether the single amplitude alone forces $\delta_F=4$ is the
question discussed in Section~\ref{sec:constrank}. For general (not necessarily quadratic)
plateaued functions, related relations between amplitudes and differential uniformity are
studied in \cite{KP} (e.g.\ \cite[Theorem~5.6]{KP}); we do not need them.
\end{remark}

\begin{remark}
The permutation hypothesis is needed: $x^5$ on $\F_{2^8}$ is quadratic with
$\delta=4$ but has components of ranks $4$ and $8$.
\end{remark}

\begin{proposition}[Isotropic spread]\label{prop:spread}
If $s_a=2$ for all $a\neq 0$, then $\Sigma=\{K_a:a\neq 0\}$ is a line spread of $V$
(a partition of $V\setminus\{0\}$ into 2-dimensional subspaces minus $0$), and
$\beta(L,L)=0$ for every $L\in\Sigma$.
\end{proposition}
\begin{proof}
If $u\in K_a\setminus\{0,a\}$ then $a\in K_u$ by symmetry of $\beta$, and $u\in K_u$, so
$K_u\supseteq\langle a,u\rangle=K_a$ and equality holds by dimension. Hence every
$x\neq 0$ lies in $K_x$, and if $x\neq 0$ lies in $K_a\cap K_b$ then $K_a=K_x=K_b$: the
$K_a$ partition $V\setminus\{0\}$. Finally, for $L=K_a$ and $y\in L\setminus\{0\}$ we have
$K_y=L$, so $\beta(y,L)=0$; thus $\beta(L,L)=0$.
\end{proof}
\begin{remark}
The isotropy of the kernels (elements of $K_a$ are pairwise $\beta$-orthogonal) appears in
\cite[Lemma~1]{MTX19} for quadratic functions with $\delta_F=4$; what is added here is the
spread structure and its consequences below.
\end{remark}

\begin{corollary}[Spread correspondence]\label{cor:spread-map}
Under the hypotheses of Theorem~\ref{thm:constrank}, the restriction $F|_L$ is $\F_2$-linear for
every $L\in\Sigma$, and $\{F(L):L\in\Sigma\}$ is a line spread of the codomain.
\end{corollary}
\begin{proof}
$\beta$ vanishes on $L\times L$ and $F(0)=0$, so $F|_L$ is additive; injectivity of $F$
makes $F(L)$ a 2-dimensional subspace and distinct images meet only in $0$. The images
cover the codomain because $F$ is bijective with $F(0)=0$.
\end{proof}

\begin{lemma}[Partial normal form]\label{lem:normal}
Under the hypotheses of Theorem~\ref{thm:constrank} there are bases of the domain and the
codomain such that $\beta(e_{2i},e_{2i+1})=0$ for $0\le i<n/2$ and
$\beta(e_0,e_j)=e_{j-2}$ for $2\le j\le n-1$.
\end{lemma}
\begin{proof}
Spread lines $L_1,\dots,L_{n/2}$ with $V=L_1\oplus\cdots\oplus L_{n/2}$ exist. Let $U$ be a
direct sum of $k/2$ spread lines ($k\le n-2$, $\dim U=k$). A spread line either lies in $U$,
or meets $U$ in exactly one point and has two points outside $U$, or is skew to $U$; and
distinct spread lines meet $U$ in distinct points. If no spread line were skew to $U$, every
one of the $2^n-2^k$ points outside $U$ would lie on a spread line of the second kind, and
there are at most $2^k-1$ such lines, carrying at most $2(2^k-1)<2^n-2^k$ outside points. So
some spread line is skew to $U$, and induction gives the decomposition. Take $e_{2i},e_{2i+1}$ a
basis of $L_{i+1}$. As $K_{e_0}=L_1=\langle e_0,e_1\rangle$, the linear map $\beta(e_0,\cdot)$ has
kernel $\langle e_0,e_1\rangle$ and rank $n-2$, so the vectors $\beta(e_0,e_j)$, $2\le j\le n-1$,
are linearly independent; compose $F$ with a linear permutation of the codomain mapping them to
$e_0,\dots,e_{n-3}$.
\end{proof}

\section{Desarguesian kernel spreads}\label{sec:f4}
A quadratic map is \emph{$\F_{2^t}$-quadratic} if its polar map $\beta$ is $\F_{2^t}$-bilinear for
some structures of $\F_{2^t}$-vector space on the domain and the codomain. For example, over
$\F_{2^n}$ with $t\mid n$ every map $\sum_{i<j}c_{ij}x^{2^{ti}+2^{tj}}+\ell(x)$ with
$c_{ij}\in\F_{2^n}$ and $\ell$ an $\F_2$-linear map is $\F_{2^t}$-quadratic; this includes the Gold
maps $x^{2^k+1}$ with $t\mid k$ plus arbitrary linear maps. All proofs in this section rest on
the following elementary fact.
\begin{lemma}\label{lem:trace-rad}
Let $E=\F_{2^t}$, let $X$ be an $E$-vector space of dimension $m$ and let $C:X\times X\to E$ be
an alternating $E$-bilinear form. Then $B=\Tr_{E/\F_2}\circ C$ is an alternating $\F_2$-bilinear
form with $\rad B=\rad C$. In particular $\dim_{\F_2}\rad B=t(m-\rk_EC)$, where $\rk_EC$ is even.
\end{lemma}
\begin{proof}
Clearly $B$ is alternating and $\rad C\subseteq\rad B$. If $u\notin\rad C$, then $C(u,\cdot)$ is
a nonzero $E$-linear map $X\to E$, hence surjective; as $\Tr_{E/\F_2}$ is not identically zero on
$E$, there is $y$ with $B(u,y)=\Tr_{E/\F_2}(C(u,y))\neq0$, so $u\notin\rad B$. Finally $\rad C$ is
an $E$-subspace of $E$-dimension $m-\rk_EC$, and alternating forms have even rank.
\end{proof}

\begin{theorem}\label{thm:general-t}
Let $F$ be a quadratic permutation of $\F_2^n$ with $s_a=t$ for every $a\neq 0$ (so
$\delta_F=2^t$ and $F$ is differentially two-valued). Suppose that the kernels form a
Desarguesian $t$-spread, that is, there is an $\F_{2^t}$-vector space structure on $\F_2^n$ (in
particular $t\mid n$) whose one-dimensional subspaces are the kernels: $K_x=\F_{2^t}x$ for every
$x\neq 0$. Then $n/t$ is odd.
\end{theorem}
\begin{proof}
Write $\Tr=\Tr_{\F_{2^t}/\F_2}$ and $m=n/t$. Since $\lambda x\in K_x$ for $\lambda\in\F_{2^t}$,
$\beta(x,\lambda x)=0$; polarising gives $\beta(\lambda x,y)=\beta(x,\lambda y)$. Let
$\{\theta_j\}$ be a basis of $\F_{2^t}/\F_2$ and $\{\theta_j^*\}$ its trace-dual basis, and put
$\Phi(x,y)=\sum_j\beta(\theta_jx,y)\otimes\theta_j^*\in\F_2^n\otimes_{\F_2}\F_{2^t}$. Expanding
$\theta_j\lambda=\sum_k\Tr(\theta_j\lambda\theta_k^*)\theta_k$ gives
\[
\Phi(\lambda x,y)=\sum_k\beta(\theta_kx,y)\otimes\sum_j\Tr(\theta_j\lambda\theta_k^*)\theta_j^*
=\sum_k\beta(\theta_kx,y)\otimes\lambda\theta_k^*=\lambda\,\Phi(x,y),
\]
and $\Phi(x,\lambda y)=\Phi(\lambda x,y)$ by self-adjointness. $\Phi$ is symmetric (by
self-adjointness and symmetry of $\beta$), $\Phi(x,x)=0$, and
$(\mathrm{id}\otimes\Tr)\Phi(x,y)=\beta\bigl((\sum_j\Tr(\theta_j^*)\theta_j)x,y\bigr)=\beta(x,y)$.
Hence $B_b=\Tr\circ C_b$ with $C_b=(\langle b,\cdot\rangle\otimes\mathrm{id})\circ\Phi$ an
alternating $\F_{2^t}$-form on $\F_{2^t}^m$, and by Lemma~\ref{lem:trace-rad}
$\dim_{\F_2}R_b=t(m-\rk_{\F_{2^t}}C_b)$ with $\rk_{\F_{2^t}}C_b$ even. If $m$ is even these
dimensions are multiples of $2t$;
they are nonzero because a component with $R_b=0$ would be bent, hence not balanced. Then
Lemma~\ref{lem:count} gives $(2^n-1)(2^t-1)=\sum_b(2^{\dim R_b}-1)\ge(2^n-1)(2^{2t}-1)$, a
contradiction.
\end{proof}
\begin{remark}
For $t\ge 3$ the argument of Proposition~\ref{prop:spread}, which uses $t=2$, does not show that
the kernels form a spread when $s_a=t$ for all $a$; this is why the spread structure is part of
the hypothesis.
\end{remark}
\begin{corollary}\label{cor:general-t}
\begin{enumerate}
\item[(i)] Every quadratic APN permutation of $\F_2^n$ has $n$ odd.
\item[(ii)] Let $F=\sum\nu_rx^{2^i+2^j}$ be a Dembowski--Ostrom polynomial over $\F_{2^n}$ all of
whose exponents have $t\mid i$ and $t\mid j$, where $t\mid n$. If $F$ is a differentially
two-valued $\{0,2^t\}$ permutation, then $n/t$ is odd.
\end{enumerate}
\end{corollary}
\begin{proof}
(i) is the case $t=1$: the kernel spread is the set of points. (ii) Here $\beta$ is
$\F_{2^t}$-bilinear, so $\F_{2^t}a\subseteq K_a$, with equality since $s_a=t$.
\end{proof}
\begin{remark}
Part (i) is classical, see e.g.\ \cite{Carlet21}; the case $t=2$ of Theorem~\ref{thm:general-t}
is Corollary~\ref{thm:desarg} below. The binomials of \cite[Theorem~1.2]{BTT12} over
$\F_{2^{3k}}$ (see also \cite[Theorem~9]{CP19}) are of the form in (ii), which explains the
hypothesis there that $k/t$ is odd; for that family its necessity is proved directly in
\cite[Lemma~2.1]{BTT12}, where $F(\gamma x)=F(x)$ for an element $\gamma$ of order $2^t+1$ when
$k/t$ is even. The quadrinomials of \cite[Problem~5.1]{BTT12} carry the same hypothesis; Qu,
Xiong and Li \cite{QXL13} showed, with the help of a computer, that they are differentially
$4$-uniform but need not be permutations. The
condition of Theorem~\ref{thm:general-t} cannot be weakened: the Gold permutations $x^{2^k+1}$
with $\gcd(k,n)=t$ and $n/t$ odd have Desarguesian kernel spreads.
\end{remark}

\begin{corollary}[Desarguesian obstruction]\label{thm:desarg}
Let $n\equiv 0\pmod 4$ and let $F$ be a quadratic permutation of $\F_2^n$ with
$\delta_F=4$. Then the isotropic spread $\Sigma$ of Proposition~\ref{prop:spread} is not
Desarguesian.
\end{corollary}
\begin{proof}
By Theorem~\ref{thm:constrank}, $s_a=2$ for all $a\neq0$. If $\Sigma$ were Desarguesian,
Theorem~\ref{thm:general-t} with $t=2$ would give that $n/2$ is odd.
\end{proof}
\begin{remark}\label{rem:F4form}
For $t=2$ the form $\Phi$ of the proof of Theorem~\ref{thm:general-t} is explicit; we state this
in the generality needed in Proposition~\ref{prop:local}. Let $X$ be an $\F_2$-space with an
$\F_4$-structure given by $\omega\in GL(X)$ with $\omega^2+\omega+1=0$, and let $w\in\F_4$
with $w^2=w+1$; the trace-dual basis of $\{1,w\}$ is $\{w^2,1\}$. Let $B$ be any
alternating $\F_2$-form on $X$ with $B(x,\omega x)=0$ for all $x$. Polarising gives
$B(\omega x,y)=B(x,\omega y)$, and
\[
C(x,y)=B(\omega x,y)+w^2B(x,y)\in\F_4
\]
satisfies $C(\omega x,y)=B(\omega^2x,y)+w^2B(\omega x,y)=wB(\omega x,y)+B(x,y)=w\,C(x,y)$ and
$C(x,\omega y)=C(\omega x,y)$; it is symmetric with $C(x,x)=0$, hence an alternating
$\F_4$-bilinear form, and $\Tr_{\F_4/\F_2}\circ C=B$ since $\Tr(1)=0$ and $\Tr(w^2)=1$. By
Lemma~\ref{lem:trace-rad}, $\rk_{\F_2}B=2\rk_{\F_4}C\equiv0\pmod4$.
\end{remark}

\begin{corollary}\label{thm:f4}
Let $n=2m$ with $m$ even and let $F$ be an $\F_4$-quadratic map on $\F_2^n$. Then no component of
$F$ has rank $n-2$. In particular no $\F_4$-quadratic map, for instance no map
$x^{2^k+1}+\ell(x)$ on $\F_{2^n}$ with $k$ even and $\ell$ linear, is a permutation with
$\delta_F=4$.
\end{corollary}
\begin{proof}
Every $\F_2$-linear functional on an $\F_4$-space is $\Tr_{\F_4/\F_2}\circ\lambda$ for an
$\F_4$-linear functional $\lambda$. Hence $B_b=\Tr_{\F_4/\F_2}\circ C_b$ with $C_b=\lambda_b\circ\beta$
alternating and $\F_4$-bilinear, and by Lemma~\ref{lem:trace-rad}
$\rk B_b=2\rk_{\F_4}C_b\equiv0\pmod4$, whereas $n-2\equiv2\pmod4$. The last statement follows
from Theorem~\ref{thm:constrank}.
\end{proof}
For power maps alone the last statement is classical: $x^{2^k+1}$ is a permutation of
$\F_{2^n}$ only if $n/\gcd(k,n)$ is odd, which for $4\mid n$ forces $4\mid\gcd(k,n)$ and hence
$\delta=2^{\gcd(k,n)}\ge16$ \cite{Gold68,Nyb93} (see also \cite{ACPBSN25}). What is new is that
adding a linear map, or passing to arbitrary $\F_4$-quadratic maps or to Desarguesian kernel
spreads, does not help.
\begin{remark}
For $m$ odd the obstruction disappears and $x^{5}$ on $\F_{2^{2m}}$ is an example. Two results
in the literature are close. For $\F_q$-quadratic maps, \cite[Theorem~5]{MTX19} proves
$\delta_F\ge q$ (and computes the boomerang uniformity when equality holds for a permutation).
For Dembowski--Ostrom polynomials $\sum\nu_r x^{2^i+2^j}$ with $t\mid i$, $t\mid j$, $t\mid n$
and without bent components, \cite[Theorem~8]{CP19} shows that $F$ is differentially two-valued
$\{0,2^t\}$ if and only if every component has radical dimension exactly $t$ (their proof
observes that radicals are $\F_{2^t}$-subspaces). Neither result contains a condition on $n/t$.
The additional ingredient here is that the radical is the radical of an alternating
$\F_4$-form, so its $\F_4$-codimension is even; this forces $m$ to be odd, and
Theorem~\ref{thm:general-t} extends this from $\F_4$-bilinear $\beta$ to Desarguesian kernel
spreads. We have not found this parity condition in the literature on differentially uniform
functions or on constant-rank spaces of alternating forms (Section~\ref{sec:constrank} discusses
the latter).
\end{remark}
\begin{remark}
Every line spread of $PG(3,2)$ is regular (a regulus of $PG(3,2)$ has only three lines), hence
Desarguesian by the Bruck--Bose correspondence \cite{BB64,Hir85}; so Corollary~\ref{thm:desarg}
contains the case $n=4$ (Section~\ref{sec:constrank} gives another short argument for it). In dimension $6$ ($n\equiv 2\pmod 4$) both kinds occur: by
the classification in Section~\ref{sec:n8search}, the component spaces of quadratic DU-4
permutations of $\F_2^6$ form exactly two classes, that of $x^5$ (Desarguesian kernel spread) and
one with a non-Desarguesian kernel spread. For $n=8$ non-Desarguesian kernel spreads are
excluded as well, by Theorem~\ref{thm:n8}.
\end{remark}

\section{Non-existence for \texorpdfstring{$n=8$}{n=8}}\label{sec:n8search}
\begin{theorem}[computer-assisted]\label{thm:n8}
There is no quadratic differentially $4$-uniform permutation of $\F_2^8$. More precisely,
$\Alt(8,\F_2)$ has no $8$-dimensional subspace $W$ such that
\begin{enumerate}
\item[(i)] every nonzero element of $W$ has rank $6$, and
\item[(ii)] every $a\in\F_2^8\setminus\{0\}$ lies in the radicals of exactly three nonzero
elements of $W$.
\end{enumerate}
Consequently every quadratic map $F:\F_2^8\to\F_2^8$ all of whose derivative kernels have
dimension $2$ (equivalently, whose differential spectrum is $\{0,4\}$) has a bent component.
\end{theorem}
For a subspace $U\subseteq\Alt(n,\F_2)$ and $a\neq0$ put $P_a(U)=\{w\in U:a\in\rad w\}$, a
subspace; (ii) says $\dim P_a(W)=2$ for all $a\neq0$.

Theorem~\ref{thm:n8} does not say that $\Alt(8,\F_2)$ has no $8$-dimensional subspace of constant
rank $6$. It excludes exactly the subspaces of constant rank $6$ that also satisfy (ii), the
condition imposed by the derivative kernels of a quadratic map with differential spectrum
$\{0,4\}$. Whether (i) alone can be satisfied is Conjecture~H of Section~\ref{sec:constrank},
which remains open for $n=8$.

\paragraph{Reduction.} Let $F$ be a quadratic DU-4 permutation of $\F_2^8$ with polar map
$\beta$. By Theorem~\ref{thm:constrank} every $B_b$ ($b\neq0$) has rank $6$ and $s_a=2$ for all
$a\neq0$. The map $b\mapsto B_b$ is linear and injective, so $W=\{B_b\}$ is $8$-dimensional and
satisfies (i). For $a\neq0$ we have $B_b(a,\cdot)=0$ if and only if
$b\perp\mathrm{Im}\,\beta(a,\cdot)$, so $\dim P_a(W)=8-\rk\beta(a,\cdot)=s_a=2$, which is (ii).
For the last statement, if all $s_a=2$ and no component is bent, all ranks are at most $6$, and
Lemma~\ref{lem:count} gives $\sum_b(2^{8-\rk B_b}-1)=3\cdot255$ with every term at least $3$, so
every component has rank $6$ and $W=\{B_b\}$ again satisfies (i) and (ii).

We call $W$ as in Theorem~\ref{thm:n8} a \emph{solution}. In the rest of this section and in
Appendix~\ref{app:search}, a \emph{point} is a nonzero vector of $\F_2^8$; the letters $a$, $r$
denote points and $x,y,c,u,w$ denote forms. The theorem is proved by a search over
subspaces $U$ of a hypothetical solution, grown one dimension at a time. For a subspace $U$ all
of whose nonzero elements have rank $6$, call a point $a$ \emph{full} for $U$ if $\dim P_a(U)=2$, and
call a coset $x+U$ with $x\notin U$ \emph{admissible} if all its elements have rank $6$ and no
element of $x+U$ has a full point of $U$ in its radical. Let $\mathcal C(U)$ be the set of
admissible cosets. Counting pairs $(a,w)$ with $w\in U\setminus\{0\}$ and $a\in\rad w\setminus\{0\}$
gives $\sum_{a\neq0}(2^{\dim P_a(U)}-1)=3(2^d-1)$ for $d=\dim U$; so if $d<8$,
some point $a$ has $\dim P_a(U)\le1$.

\begin{lemma}[radical closure]\label{lem:closure}
Let $W$ be a solution, $U\subseteq W$ a subspace of dimension $d<8$ and $r\neq0$ a point with
$\dim P_r(U)\le1$. Then $W$ contains an admissible coset $c+U$ that contains an element whose
radical contains $r$. Moreover:
\begin{enumerate}
\item[(P1)] $\dim P_a(U)\le2$ for all $a\neq0$;
\item[(P2)] $|\mathcal C(U)|\ge 2^{8-d}-1$;
\item[(P3)] for every $a$ with $\dim P_a(U)=1$ (resp.\ $0$), at least one (resp.\ three)
admissible cosets contain an element whose radical contains $a$.
\end{enumerate}
\end{lemma}
\begin{proof}
Since $\dim P_r(W)=2>\dim P_r(U)$, there is $c\in P_r(W)\setminus U$. The coset $c+U$ lies in
$W\setminus\{0\}$, so its elements have rank $6$; if some $c+u$ had a full point $a$ of $U$ in
its radical, then $P_a(W)\supseteq P_a(U)+\langle c+u\rangle$ would have dimension $3$. So $c+U$
is admissible, and $r\in\rad c$. (P1) holds as $P_a(U)\subseteq P_a(W)$. The same argument shows
that all $2^{8-d}-1$ nonzero cosets of $W/U$ are admissible, which is (P2). For (P3), the image
of $P_a(W)$ in $W/U$ has dimension $2-\dim P_a(U)$; its nonzero elements are admissible cosets
containing an element of $P_a(W)$.
\end{proof}

\paragraph{The search.} All forms of rank $6$ on $\F_2^8$ are equivalent, so we may assume that
a solution contains $J=e_0\wedge e_1+e_2\wedge e_3+e_4\wedge e_5$ (radical
$\langle e_6,e_7\rangle$). Starting from $U_1=\langle J\rangle$, a node $U$ of dimension $d<8$ is
discarded if (P2) fails, or, for $d\ge3$, if (P3) fails ((P1) holds for every node that is
reached, by Lemma~\ref{lem:B-step}(b)); otherwise a point $r$ with
$\dim P_r(U)\le1$ is chosen by a fixed rule (for $d\le2$ the first point with $\dim P_r(U)=1$,
or with $\dim P_r(U)=0$ if there is none; for $d\ge3$ the point minimising the number of
candidates below), and the children of $U$ are the subspaces $U+\langle c\rangle$ for the
admissible cosets $c+U$ containing an element whose radical contains $r$ (the
\emph{candidates}), one for each orbit of a group $H\le GL(8,2)$ that stabilises $U$ and fixes
$r$. Here $g\in GL(8,2)$ acts by $B^g(v,v')=B(gv,gv')$; it preserves ranks and maps solutions to
solutions, since $\rad B^g=g^{-1}\rad B$. By Lemma~\ref{lem:closure}, every solution containing
$U$ is mapped by an element of $H$ to a solution containing one of the children, so every
solution is equivalent to one containing a node of dimension $8$; these are tested against (ii).
The search is small because each step prescribes a point in the radical of the new element
and because admissibility removes many cosets once full points appear
(Appendix~\ref{app:small}).

Appendix~\ref{app:search} states the search as a procedure, proves its completeness
(Theorem~\ref{thm:B-complete}) and describes how the programs implement each step.

\paragraph{Symmetry.} For $U_1$ we use explicit generators of subgroups of the stabiliser of $J$:
the symplectic transvections $x\mapsto x+J(v,x)v$ for $v\in\langle e_0,\dots,e_5\rangle$ of
weight $1$ or $2$, the group $GL(2,2)$ acting on $\langle e_6,e_7\rangle$, and the maps
$x\mapsto x+x_ie_6$, $x\mapsto x+x_ie_7$ ($i\le5$). Stabilisers of points and of cosets are
generated by random products corrected by an orbit transversal and by random Schreier
generators. A proper subgroup of a stabiliser only produces more children, so correctness does
not depend on these subgroups being the full stabilisers; every generator of every group used is
checked to stabilise $U$ (and to fix $r$). For the two top levels the program writes a
certificate: for each of the $740{,}736$ candidates of $U_1$ and the $727{,}936$ candidates of
the three nodes of dimension $2$, an explicit matrix of $H$ mapping an orbit representative onto
it.

\paragraph{Results.} The search has $3$ nodes of dimension $2$, $486$ of dimension $3$,
$1{,}466{,}318$ of dimension $4$ and $244{,}686{,}904$ of dimension $5$. Every node of dimension
$5$ fails (P2) ($244{,}668{,}896$ nodes) or (P3) ($18{,}008$ nodes), so no node of dimension $6$
is reached and no solution exists. This proves Theorem~\ref{thm:n8}. The run consisted of two
parallel processes, which took $2.2$ and $2.1$ hours ($4.3$ CPU-hours in total) on a two-core
virtual machine with $2.1$\,GHz Intel Xeon processors; a second run from the archived sources
gave the same numbers at every dimension.

\paragraph{Verification.}
\begin{itemize}
\item \emph{A second implementation of the same search.} A second program
(Python/\texttt{numpy}) was written separately from the mathematical specification of the
search (including its tie-breaking rule), without access to the code of the first program. It
belongs to the same study, so it is a separate code path rather than a third-party
replication. It recomputes the table
of rank-$6$ forms by Pfaffians (rank $6$ if and only if the $8\times8$ Pfaffian vanishes and one of
the $28$ principal $6\times6$ Pfaffians does not) and obtains the same $149{,}920{,}960$ forms as
the first program, which uses Gaussian elimination. It checks the certificate of the two top
levels line by line: each matrix is invertible, stabilises the node, fixes $r$ and maps the
representative onto the candidate; the listed candidates are exactly the candidates computed
from the definitions; and the $486$ nodes of dimension $3$ are exactly those of the first
program. It then searches the subtrees below the $486$ nodes of dimension $3$ without any
symmetry reduction, so that no symmetry argument below dimension $3$ is used. This search has
$8{,}604{,}140$ nodes of dimension $4$ and $1{,}381{,}518{,}812$ of dimension $5$ ($13.4$
hours in total for two parallel processes); again every node of dimension $5$ fails (P2) or (P3), and no solution exists. So the
certificate of the two top levels and this search verify Theorem~\ref{thm:n8} without using the
first program (Appendix~\ref{app:depend} lists what each part of the proof depends on). On the
subtrees of five nodes of dimension
$3$ (from all three nodes of dimension $2$), the first program run without symmetry below
dimension $3$ gives exactly the same numbers of nodes and of discarded nodes at every dimension.
Agreement of the two programs tests the implementation; the pruning rules themselves are
justified by Lemma~\ref{lem:closure}.
\item \emph{Dimension $6$.} The same program for $n=6$ (rank $4$, dimension $6$) finds exactly
two $GL(6,2)$-classes of solutions: the component space of $x^5$ (Desarguesian kernel spread,
stabiliser of order $362{,}880=|\Gamma L(3,4)|$) and a class with a non-Desarguesian kernel
spread in which $50$ of the $210$ pairs of kernel lines are closed (stabiliser of order
$5{,}760$). Both are component spaces of DU-4 permutations. An affine equivalence
$G=A_2\circ F\circ A_1$ changes the component space only through the linear part of $A_1$ (the
output map $A_2$ only re-indexes the components, and affine terms do not change $\beta$), so
each affine class of quadratic DU-4 permutations of $\F_2^6$ determines one of these two classes;
the eight affine classes with $\delta=4$ found in \cite{DMB19} are distributed over our two
classes, both of which occur (one component space carries several permutations, which differ in
their linear parts). Without symmetry reduction the program finds
$12{,}288$ solutions containing $J$; this equals
$|GL(6,2)|\cdot63\cdot(1/362{,}880+1/5{,}760)/18{,}228=192+12{,}096$, the number predicted by the
two stabilisers ($18{,}228$ is the number of forms of rank $4$ in $\Alt(6,\F_2)$). A plain
enumeration that does not prescribe radical points (a depth-first search over all admissible
cosets that reaches each subspace containing $J$ exactly once, through its greedy basis modulo
$J$) finds the same $12{,}288$ subspaces. The same enumeration with the condition on full points
dropped, so that a coset is only required to consist of forms of rank $4$, again finds exactly
these $12{,}288$ subspaces. So for $n=6$ condition (ii) follows from (i), and $\Alt(6,\F_2)$ has
exactly two $GL(6,2)$-classes of $6$-dimensional subspaces of constant rank $4$.
\end{itemize}

\section{Local structure of kernel spreads and earlier approaches}\label{sec:local}
For kernel lines $L\neq M$ let $k(L,M)$ be the number of kernel lines contained in $L\oplus M$,
and call $\{L,M\}$ \emph{closed} if $k(L,M)=5$, i.e.\ $L\oplus M$ is a union of kernel lines. A
Desarguesian kernel spread has $k(L,M)=5$ for all pairs. Put $d(L,M)=\dim\langle\beta(L,M)\rangle$;
since $\beta(L,L)=\beta(M,M)=0$, this is the dimension of $\langle\beta(U,U)\rangle$ for $U=L\oplus M$.
\begin{proposition}\label{prop:local}
Let $n\ge 6$ be even and let $\beta$ have $s_a=2$ for all $a\neq0$ and all components of rank
$n-2$ (as for every quadratic DU-4 permutation, Theorem~\ref{thm:constrank}). For kernel lines
$L\neq M$:
\begin{enumerate}
\item[(a)] $k(L,M)\in\{2,3,5\}$;
\item[(b)] $d(L,M)=4,3,2$ according as $k(L,M)=2,3,5$;
\item[(c)] if $n=8$, no pair is closed, so $k(L,M)\in\{2,3\}$.
\end{enumerate}
\end{proposition}
\begin{proof}
(a) Let $U=L\oplus M$ and suppose $U$ contains a third kernel line $N$. Three pairwise skew lines of
$PG(3,2)$ lie in a unique spread $\{L,M,N,P,Q\}$ of $PG(U)$; it is regular, so its lines are the
$\F_4$-points for some $\omega\in GL(U)$ with $\omega^2+\omega+1=0$. A line of $U$ skew to $L$, $M$
and $N$ lies in the six points of $P\cup Q$, hence contains two points of $P$ or of $Q$ and equals
$P$ or $Q$; so the kernel lines in $U$ are among $L,M,N,P,Q$. The map $q(x)=\beta(x,\omega x)$ is
quadratic on $U$ and satisfies $q(\omega x)=\beta(\omega x,x+\omega x)=q(x)$, so it takes one value
$q_X$ on the nonzero points of each $\F_4$-point $X$. Every component of $q$ is a quadratic form on
$\F_2^4$, which takes the value $1$ an even number of times; hence $\sum_{x\in U}q(x)=0$, i.e.\
$q_L+q_M+q_N+q_P+q_Q=0$. A line $\F_4x$ of $U$ is a kernel line if and only if $\beta(x,\omega x)=0$
(as $s_x=2$ and $\omega x\notin\{0,x\}$). So $q_L=q_M=q_N=0$, hence $q_P=q_Q$, and $P$ is a kernel
line if and only if $Q$ is. Thus $k\neq4$.

(b) Let $U=L\oplus M$ and let $N_0,N_2,N_4$ be the numbers of $b\neq0$ for which $B_b|_U$ has
rank $0$, $2$, $4$, so that $N_0+N_2+N_4=2^n-1$; since $B_b|_U=0$ exactly when
$b\perp\langle\beta(U,U)\rangle$, we have $N_0=2^{n-d}-1$. A point $a$ on one of the $k$ kernel
lines inside $U$ has $K_a\subseteq U$, so $\beta(a,\cdot)|_U$ has kernel $K_a$ and
$\dim\beta(a,U)=2$; every other point $a\in U$ has $K_a\cap U=\langle a\rangle$ and
$\dim\beta(a,U)=3$. Remark~\ref{rem:count-general} for $U$ gives
\begin{equation}\label{eq:local}
15N_0+3N_2=3k(2^{n-2}-1)+(15-3k)(2^{n-3}-1).
\end{equation}
If $k=5$, then $\beta(x,\omega x)=0$ on $U$, so by Remark~\ref{rem:F4form}, applied to each
$B_b|_U$, every $B_b|_U$ has rank $0$ or $4$; thus $N_2=0$, and \eqref{eq:local} gives
$N_0=2^{n-2}-1$, i.e.\ $d=2$. If $k\in\{2,3\}$, some point of $U$ has $\dim\beta(a,U)=3$, so
$d\in\{3,4\}$. If $k=3$, then $v:=q_P=q_Q\neq0$ by (a), since $P$ is not a kernel line; every
$b\perp v$ has $B_b(x,\omega x)=0$ on $U$, so $B_b|_U$ has rank $0$ or $4$ by
Remark~\ref{rem:F4form}, and therefore $N_2\le(2^n-1)-(2^{n-1}-1)=2^{n-1}$. Solving
\eqref{eq:local} for $N_2$ gives the following table.
\begin{center}
\begin{tabular}{ccccc}\toprule
$k$ & $d$ & $N_0$ & $N_2$ & $3N_0+N_2$\\\midrule
$2$ & $3$ & $2^{n-3}-1$ & $2^{n-2}$ & $3\cdot2^{n-2}-3-2^{n-3}$\\
$2$ & $4$ & $2^{n-4}-1$ & $9\cdot2^{n-4}$ & $3\cdot2^{n-2}-3$\\
$3$ & $3$ & $2^{n-3}-1$ & $3\cdot2^{n-3}$ & $3\cdot2^{n-2}-3$\\
$3$ & $4$ & $2^{n-4}-1$ & $11\cdot2^{n-4}>2^{n-1}$ & excluded\\
$5$ & $2$ & $2^{n-2}-1$ & $0$ & $3\cdot2^{n-2}-3$\\\bottomrule
\end{tabular}
\end{center}
So every pair has $3N_0+N_2\le3\cdot2^{n-2}-3$, with equality unless $(k,d)=(2,3)$. We show that
the average of $3N_0+N_2$ over all unordered pairs of kernel lines equals $3\cdot2^{n-2}-3$;
then no pair has $(k,d)=(2,3)$, which proves (b).

Fix $b\neq0$. The ordered pairs $(x,y)\in V^2$ with $B_b(x,y)=0$ number
$4\cdot2^n+(2^n-4)2^{n-1}$: for the $4$ vectors $x\in R_b$ every $y$ qualifies, and for the other
$x$ the functional $B_b(x,\cdot)$ is nonzero and vanishes on $2^{n-1}$ vectors. Of these pairs,
$2^{n+1}-1$ have $x=0$ or $y=0$, and $3(2^n-1)$ consist of nonzero vectors in a common kernel
line ($B_b$ vanishes on $L'\times L'$ for every kernel line $L'$, and there are
$|\Sigma|=(2^n-1)/3$ kernel lines with $9$ such ordered pairs each). This leaves
\[
Z=2^{2n-1}-3\cdot2^n+4
\]
ordered pairs of nonzero vectors in distinct kernel lines. Group them by the ordered pair
$(L',M')$ of lines. Let $\rho\in\{0,1,2\}$ be the rank of the $2\times2$ matrix of $B_b$ on
$L'\times M'$; the number of $(x,y)\in(L'\setminus0)\times(M'\setminus0)$ with $B_b(x,y)=0$ is
$9$, $5$, $3$ for $\rho=0,1,2$ (for $\rho=1$, $B_b(x,y)=\lambda(x)\mu(y)$ with nonzero functionals
$\lambda,\mu$, and there are $1\cdot3+3\cdot1-1$ zeros; for $\rho=2$ every $x$ has exactly one
partner $y$). With $n_\rho(b)$ the number of ordered pairs of distinct kernel lines of type $\rho$, we
get $n_0+n_1+n_2=|\Sigma|(|\Sigma|-1)$ and $9n_0+5n_1+3n_2=Z$, hence
\[
3n_0(b)+n_1(b)=\tfrac12\bigl(Z-3|\Sigma|(|\Sigma|-1)\bigr)=\tfrac16(2^n-4)(2^{n-1}-2).
\]
Since $B_b$ vanishes on $L'\times L'$ and on $M'\times M'$, its rank on $U'=L'\oplus M'$ is $2\rho$;
so type $0$ means $B_b|_{U'}=0$ and type $1$ means that $B_b|_{U'}$ has rank $2$. Each unordered
pair $\{L',M'\}$ gives two ordered pairs of the same type, so summing over $b\neq0$ and
exchanging the order of summation,
\[
\sum_{b\neq0}\bigl(3n_0(b)+n_1(b)\bigr)=2\sum_{\{L',M'\}}\bigl(3N_0(U')+N_2(U')\bigr).
\]
The left side is $(2^n-1)(2^n-4)(2^{n-1}-2)/6$ and there are $\binom{|\Sigma|}{2}=(2^n-1)(2^n-4)/18$
unordered pairs, so the average of $3N_0+N_2$ is
$\frac{(2^n-1)(2^n-4)(2^{n-1}-2)/12}{(2^n-1)(2^n-4)/18}=\frac32(2^{n-1}-2)=3\cdot2^{n-2}-3$.

(c) Let $n=8$ and $k=5$. By (b), $S=\langle\beta(U,U)\rangle$ has dimension $2$, so there are $63$
values $b\neq0$ with $b\perp S$, i.e.\ with $B_b|_U=0$. For each of them $U\cap R_b\neq0$:
otherwise $U\oplus R_b$ would be a totally isotropic subspace of dimension $6$ for $B_b$, whereas
a totally isotropic subspace for a form of rank $6$ on $\F_2^8$ has dimension at most $2+3=5$
(the radical has dimension $2$, and a totally isotropic subspace of the nondegenerate quotient
of dimension $6$ has dimension at most $3$). So
there are at least $63$ pairs $(a,b)$ with $a\in U\setminus\{0\}$, $b\perp S$, $b\ne0$ and $a\in R_b$.
But every $a\neq0$ lies in exactly three radicals (those of the $b\in(\mathrm{Im}\,\beta(a,\cdot))^\perp$),
so there are at most $3\cdot15=45$ such pairs, a contradiction.
\end{proof}
\begin{remark}
For $n=6$ the argument of (c) forces $R_b\subseteq U$ for the $15$ values of $b$, i.e.\ exactly $45$
pairs, and closed pairs do occur: by the classification for $n=6$ in Section~\ref{sec:n8search},
$50$ or all $210$ of the pairs are closed, the other pairs having $k=2$ and $d=4$. Part (c) gives a
second proof of Corollary~\ref{thm:desarg} for $n=8$ (it shares the $\F_4$-trace step of
Remark~\ref{rem:F4form}). For
$n=8$ the hypotheses of Proposition~\ref{prop:local} cannot be satisfied, by
Theorem~\ref{thm:n8}; parts (a) and (b) hold in every even dimension $n\ge6$.
\end{remark}


\paragraph{Earlier approaches for $n=8$.} Before the search of Section~\ref{sec:n8search} we
tried three other approaches; they agree with Theorem~\ref{thm:n8} but did not settle the case.
Programs, logs and full proofs are in the repository.
\begin{itemize}
\item \emph{Structured classes.} Representing $\F_{2^8}$ with the modulus $X^8+X^4+X^3+X^2+1$ and
using exponents $2^i+2^j$, no binomial $x^{d_1}+c\,x^{d_2}$ ($96{,}390$ polynomials), no
polynomial with coefficients in $\F_2$ ($2^{28}$ polynomials) and no trinomial
($213{,}021{,}900$ polynomials) has all $255$ components of rank $6$; the closest are $12{,}288$
polynomials with coefficients in $\F_2$ that have three components of other ranks.
\item \emph{SAT.} A direct encoding of quadratic DU-4 permutations (variables $\beta(e_i,e_j)$
and $F(e_i)$, symmetry breaking by Lemma~\ref{lem:normal}) is unsatisfiable for $n=4$ (with a
DRAT proof checked by \texttt{drat-trim} \cite{WHH14}) and satisfiable for $n=6$, but
\texttt{kissat} \cite{Kissat} gave no verdict for $n=8$ within several hours, also after
splitting by the local structure of Proposition~\ref{prop:local}. With the Desarguesian condition
added, the instance is unsatisfiable within a second (a DRAT-checked consistency check of
Corollary~\ref{thm:desarg}).
\item \emph{Symmetries.} Call $(A,B)\neq(I,I)$ in $GL(8,2)^2$ a symmetry of $\beta$ if
$\beta(Ax,Ay)=B\beta(x,y)$; every nontrivial affine self-equivalence of $F$ gives one. By a case
analysis over the conjugacy classes of the possible pairs $(A,B)$ with $A$ of prime order, using
SAT calls certified by DRAT proofs and exhaustive enumerations, we had excluded every $\beta$
with a symmetry (the proof and the computations are in the repository). This statement is now
contained in Theorem~\ref{thm:n8}.
\end{itemize}

\section{Constant-rank subspaces and an open problem}\label{sec:constrank}
For $n\ge 4$, Theorem~\ref{thm:constrank} says that the component forms of a quadratic DU-4
permutation span an $n$-dimensional subspace $W$ of $\Alt(n,\F_2)$ all of whose nonzero
elements have rank $n-2$ (the map $b\mapsto B_b$ is injective, since a nonzero $b$ gives a
component of rank $n-2>0$). For $n=4$ no such subspace exists: the forms of rank $\le 2$ in
$\Alt(4,\F_2)$ form the Klein quadric $Q^+(5,2)$, whose largest totally singular subspaces
have dimension $3$ (equivalently, the MacWilliams transform in the alternating-forms scheme
\cite{DG75} has a negative coefficient). For $n=8$ the transform is non-negative and integral,
so it gives no obstruction (Table~\ref{tab:macw}).

Constant-rank subspaces of alternating forms have mostly been studied over larger fields.
Gow and Quinlan \cite{GQ09} construct them from Galois extensions; as far as we can see, these
constructions do not give $n$-dimensional subspaces of constant rank $n-2$ when $4\mid n$ (for
instance the component space of $x^5$ on $\F_{2^8}$ contains forms of ranks $4$ and $8$). The
extension of these constructions by Gupta and Mandal \cite{GM23} assumes that the characteristic
is not two. Gow
\cite{Gow18} shows that if $q\ge m+1$, the distinct radicals of the nonzero elements of an
$n$-dimensional subspace of $\Alt(n,\F_q)$ of constant rank $m$ form a spread
\cite[Theorem~13]{Gow18}, and he bounds the dimension of constant-rank subspaces under the same
field-size hypothesis \cite[Theorem~11]{Gow18}. These results do not apply here, since $q=2$ and
$m=n-2$, and over $\F_2$ the spread conclusion indeed fails: for $n=6$, the non-Desarguesian
class of Section~\ref{sec:n8search} is a $6$-dimensional subspace of constant rank $4$ whose $63$
nonzero elements have $61$ distinct radicals, so the radicals do not form a spread (for $x^5$
they are the $21$ lines of the kernel spread, each occurring three times). In our setting the
spread comes from the derivative kernels, that is, from the condition $s_a=2$, and not from the
radicals. We are not aware of results settling the existence question over $\F_2$.

In Table~\ref{tab:macw}, $W^\perp$ is taken with respect to the pairing $\langle A,A'\rangle=\sum_{i<j}A_{ij}A'_{ij}$.
A rank-$2$ element of $\Alt(n,\F_2)$ is $u\wedge v$ (the form
$(x,y)\mapsto\langle u,x\rangle\langle v,y\rangle+\langle v,x\rangle\langle u,y\rangle$) for a
unique line $\langle u,v\rangle$, and $\langle A,u\wedge v\rangle=A(u,v)$. Hence the rank-$2$
elements of $W^\perp$ correspond to the lines $\langle u,v\rangle$ with $B(u,v)=0$ for all
$B\in W$, i.e.\ with $\beta(u,v)=0$. Counting ordered pairs $(u,v)$ with $v\in
K_u\setminus\{0,u\}$ and using Lemma~\ref{lem:count}, the number of these lines is
$\sum_a(2^{s_a}-2)/6=(2^n-1)/3$ for every $n$-dimensional subspace of constant rank $n-2$, in
agreement with the
entry $A'_1$ of Table~\ref{tab:macw}. These lines form a spread if and only if $s_a=2$ for all
$a$. This holds for DU-4 permutations (Theorem~\ref{thm:constrank}); we do not know whether it
follows from the constant-rank condition in general. For $n=6$ it does, by the enumeration at
the end of Section~\ref{sec:n8search}.
\begin{table}[h]\centering
\begin{tabular}{cccl}\toprule
$n$ & $d$ & verdict & dual rank distribution $A'_0,\dots,A'_{n/2}$\\\midrule
4 & 4 & impossible & $1,\ 5,\ -2$\\
6 & 6 & consistent & $1,\ 21,\ 210,\ 280$\\
8 & 8 & consistent & $1,\ 85,\ 20706,\ 590920,\ 436864$\\\bottomrule
\end{tabular}
\caption{MacWilliams transform of the rank distribution $(1,0,\dots,0,2^d-1,0)$ of a
$d$-dimensional subspace of $\Alt(n,\F_2)$ of constant rank $n-2$; $A'_r$ is the number of
elements of rank $2r$ in $W^\perp$.}
\label{tab:macw}
\end{table}
\begin{conjecture}[H]
For $n\equiv 0\pmod 4$, $\Alt(n,\F_2)$ has no $n$-dimensional subspace whose nonzero
elements all have rank $n-2$.
\end{conjecture}
Conjecture~H would imply that no quadratic DU-4 permutation exists when $4\mid n$, but it is a
stronger statement: a constant-rank subspace is only a necessary ingredient. A permutation
additionally needs $s_a=2$ for all $a$ and, for a quadratic map $Q$ with polar map $\beta$, a
linear map $\Lambda$ with $Q(a)+\Lambda(a)\notin\mathrm{Im}\,\beta(a,\cdot)$ for all $a\neq 0$. A
counterexample to Conjecture~H would therefore not by itself produce a permutation.
For $n=8$, Theorem~\ref{thm:n8} proves Conjecture~H for subspaces with $s_a=2$ for all $a$
(in the notation of Section~\ref{sec:n8search}, $\dim P_a(W)=2$). For a general $8$-dimensional
subspace of constant rank $6$ the counting argument above gives
$\sum_{a\neq0}(2^{\dim P_a(W)}-1)=3\cdot255$ with $\dim P_a(W)\ge1$, so either all
$\dim P_a(W)=2$, which is excluded, or some point lies in the radicals of at least seven
elements; the second case is not covered by our search, and Conjecture~H remains open for
$n=8$. (For $n=6$ the analogous second case does not occur, as noted above.)

\section{Conclusion}
By \cite{CP19}, a quadratic differentially 4-uniform permutation of $\F_2^n$ has all components
of rank $n-2$ and all derivative kernels of dimension $2$. We observed that these kernels form a
line spread that is totally isotropic for every component, and showed that when $4\mid n$ this
spread cannot be Desarguesian; the same argument shows in general that a Desarguesian kernel
$t$-spread forces $n/t$ to be odd. This explains why no $\F_4$-quadratic construction, for
instance a Gold map with an added linear part, works when $4\mid n$, gives a uniform reason for the parity
hypothesis in the binomial family of \cite{BTT12}, and places the classical non-existence of
quadratic APN permutations in even dimension in the same framework. For $n=8$ we showed, by an
exhaustive search over the possible spaces of components, that no quadratic DU-4 permutation
exists (DU-4 permutations of higher degree do exist, the inverse map being the best-known
example); in fact the search never gets beyond five of the eight dimensions. Together with $n=4$,
this suggests that quadratic DU-4 permutations of $\F_2^n$, $n$ even, exist only for
$n\equiv2\pmod4$, where the Gold maps provide them for every $n\ge6$. The next open case is $n=12$, which is far out of reach of the same search. Theorem~\ref{thm:n8}
excludes only the constant-rank subspaces that satisfy the radical condition (ii);
Conjecture~H, which drops this condition, is open already for $n=8$.

\section*{Reproducibility}
The programs, logs and certificates are available in the repository
\begin{center}\url{https://github.com/AlexGuseinov/quadratic-perm-4n}.\end{center}
Its file \texttt{ARTIFACT.md} lists, for each computational statement of this paper, the command
that reproduces it, the expected output and what has to be trusted;
\texttt{expected\_counts.json} contains the numbers of nodes of both programs in
machine-readable form, and the checking scripts compare their logs with it.

\emph{First program} (directory \texttt{search\_n8/}). The search of
Section~\ref{sec:n8search} is \texttt{crsub\_enum.c} (C,
compiled with \texttt{gcc 13.3 -O3 -march=native}; it needs an x86-64 processor with the BMI2
instructions). The script \texttt{run\_n8.sh} runs it in two parallel parts and records the
compiler, the command lines and the SHA-256 hashes of the sources, the binary and the table of
rank-$6$ forms; \texttt{check\_n8\_logs.py} checks that all $486$ nodes of dimension $3$ were
processed, that no node of dimension $6$ occurred and that the numbers of nodes are the expected
ones. All random choices in the computation of stabilisers are derived from a fixed seed and the
path from the root, so the program is deterministic and a rerun reproduces the same tree and the
same numbers. The reported run used Ubuntu 24.04 on a two-core virtual machine with Intel Xeon
processors at $2.1$\,GHz: two parallel processes of $2.2$ and $2.1$ hours ($4.3$ CPU-hours), at
most $0.3$\,GB of memory per process. The second run, from the archived sources, took $3.4$ and
$3.3$ hours because it shared the machine with the second program.

\emph{Second program.} The directory \texttt{second\_impl/} contains the second implementation
(Python 3.11 with NumPy 2.4), the certificate of the two top levels (\texttt{cert8.txt.xz},
$1{,}468{,}672$ lines, with its SHA-256 hash), its checker and the logs of the complete run over
the $486$ nodes of dimension $3$ (two parallel processes of $6.8$ and $6.7$ hours, $13.4$ hours
in total, partly on the shared machine; at most $0.6$\,GB of memory per process). One command,
\texttt{sh second\_impl/run\_all.sh}, carries out the verification of Theorem~\ref{thm:n8} that
does not use the first program: it computes the table of rank-$6$ forms by Pfaffians, checks the
certificate, searches the $486$ subtrees and compares the result with the expected numbers. The
only data it takes from the first program are the certificate and the list of the $486$ nodes of
dimension $3$, and it checks both against the definitions; the checker shares no code with the
program that wrote the certificate.

\emph{Small computations.} The script \texttt{check\_all.sh} re-checks, in about a minute, the
small computations only (rank counts, the MacWilliams table, the SAT model for $n=4$ and $n=6$,
the Desarguesian instance for $n=8$, and the examples for $x^5$); it does not run the search for
$n=8$. The enumerations for $n=6$ of Section~\ref{sec:n8search}, the searches over structured
classes, the symmetric-case computations and the long solver runs mentioned in
Section~\ref{sec:local} are documented in the repository by their commands, logs and
certificates.


\appendix
\section{The computation in mathematical terms}\label{app:search}
This appendix describes the search of Section~\ref{sec:n8search} in mathematical terms: the
objects it works with, the procedure and a proof that it is complete, how each quantity is
computed, how the symmetry reduction is certified, and what each part of the proof of
Theorem~\ref{thm:n8} depends on. Throughout $n=8$, $V=\F_2^8$, and forms are elements of
$\Alt(8,\F_2)$.

\subsection{The objects}
A form $B$ is determined by its $28$ values $B(e_i,e_j)$, $i<j$; this identifies
$\Alt(8,\F_2)$ with $\F_2^{28}$, and subspaces of forms with subspaces of $\F_2^{28}$. In the
programs and their logs the pairs are ordered lexicographically, $(0,1),(0,2),\dots,(6,7)$, a
form is written as the integer whose $k$-th binary digit is its $k$-th coordinate, and a point
$a$ as the integer $\sum_ia_i2^i$; thus $J=2^0+2^{13}+2^{22}=4{,}202{,}497$. A table
of all $2^{28}$ forms records which of them have rank $6$. The first program computes ranks by
Gaussian elimination; the second uses Pfaffians: $\rk B=8$ if and only if the Pfaffian of $B$ is
nonzero, and $\rk B\ge6$ if and only if one of the $28$ principal $6\times6$ Pfaffians is
nonzero. Both give the same $149{,}920{,}960$ forms of rank $6$, the number given by the classical
formula for the number of alternating forms of given rank (see \cite{DG75}).

If the basis of a subspace $U$ is in reduced echelon form (each basis vector has a pivot
coordinate in which the other basis vectors vanish), every coset $y+U$ contains exactly one
vector that vanishes in all pivot coordinates; it is obtained from $y$ by adding the basis
vectors whose pivot coordinate is $1$ in $y$. Sets of cosets are stored as sets of these
representatives.

For $a\in V$ let $\varphi_a:\Alt(8,\F_2)\to V^*$, $\varphi_a(B)=B(a,\cdot)$, that is,
$\varphi_a(B)(e_k)=\sum_ia_iB(e_i,e_k)$; it is linear (an $8\times28$ matrix over $\F_2$), and
$a\in\rad B$ if and only if $\varphi_a(B)=0$. For a subspace $U$ we have
$P_a(U)=U\cap\ker\varphi_a$, so
\[
\dim P_a(U)=\dim U-\dim\varphi_a(U).
\]
A \emph{node} is a subspace $U$ all of whose nonzero elements have rank $6$ and with
$\dim P_a(U)\le2$ for all $a\neq0$. For a node $U$ let $\mathcal F(U)=\{a\neq0:\dim P_a(U)=2\}$ (the
full points). A coset $x+U$, $x\notin U$, is admissible if all its elements have rank $6$ and
$\varphi_a(x)\notin\varphi_a(U)$ for every $a\in\mathcal F(U)$ (by the display, this says that no
element of $x+U$ has $a$ in its radical); $\mathcal C(U)$ is the set of admissible cosets. For
a point $r$ put
$\mathrm{Cand}(U,r)=\{x+U\in\mathcal C(U):\varphi_r(x)\in\varphi_r(U)\}$.
For $g\in GL(V)$ and a form $B$ put $B^g(v,v')=B(gv,gv')$ for $v,v'\in V$, and
$U^g=\{u^g:u\in U\}$.


\subsection{Two lemmas}
\begin{lemma}[group action]\label{lem:B-action}
Let $g,h\in GL(V)$, $B$ a form, $U$ a subspace and $a\in V$.
\begin{enumerate}
\item[(a)] $B\mapsto B^g$ is linear and bijective, $(B^g)^h=B^{gh}$ and $\rk B^g=\rk B$.
\item[(b)] $a\in\rad B^g$ if and only if $ga\in\rad B$; hence $P_a(U^g)=P_{ga}(U)^g$.
\item[(c)] $U^g$ is a node if and only if $U$ is, and then $\mathcal F(U^g)=g^{-1}\mathcal F(U)$.
If $W$ is a solution, so is $W^g$.
\item[(d)] If $U$ is a node, $U^g=U$ and $gr=r$, then $x+U\mapsto x^g+U$ is a permutation of
$\mathcal C(U)$ and of $\mathrm{Cand}(U,r)$.
\end{enumerate}
\end{lemma}
\begin{proof}
(a) In matrix form $B^g=g^{\mathsf T}Bg$. (b) $B^g(a,v)=B(ga,gv)$ vanishes for all $v\in V$ if and
only if $B(ga,\cdot)=0$. Hence $w^g\in P_a(U^g)$ if and only if $w\in P_{ga}(U)$. (c) follows from
(a) and (b), since $\dim P_a(U^g)=\dim P_{ga}(U)$. (d) By (c), $g\mathcal F(U)=\mathcal F(U)$.
For $y\in x+U$, $y^g$ has rank $6$ if and only if $y$ has, and by (b) $y^g$ has a point $a$ of
$\mathcal F(U)$ in its radical if and only if $y$ has the point $ga\in\mathcal F(U)$ in its
radical; likewise $r\in\rad y^g$ if and only if $r=gr\in\rad y$. So $x+U$ and $x^g+U=(x+U)^g$
are admissible (resp.\ candidates) together; the inverse map is given by $g^{-1}$.
\end{proof}

\begin{lemma}[one step]\label{lem:B-step}
Let $U$ be a node of dimension $d$, $x+U\in\mathcal C(U)$ and $U'=U+\langle x\rangle$. For
$a\neq0$ put $\varepsilon_a=1$ if $\varphi_a(x)\in\varphi_a(U)$ and $\varepsilon_a=0$ otherwise.
\begin{enumerate}
\item[(a)] $\dim P_a(U')=\dim P_a(U)+\varepsilon_a$ for all $a\neq0$.
\item[(b)] $\varepsilon_a=0$ for $a\in\mathcal F(U)$. Hence $U'$ is a node and
$\mathcal F(U)\subseteq\mathcal F(U')$.
\item[(c)] For $y\notin U'$, the coset $y+U'$ is admissible for $U'$ if and only if $y+U$ and
$y+x+U$ are admissible for $U$ and $\varphi_a(y)\notin\varphi_a(U')$ for all
$a\in\mathcal F(U')\setminus\mathcal F(U)$.
\item[(d)] $x+U$ contains an element with $r$ in its radical if and only if
$\varphi_r(x)\in\varphi_r(U)$.
\item[(e)] If $d<8$, some $a$ has $\dim P_a(U)\le1$.
\end{enumerate}
\end{lemma}
\begin{proof}
(a) $\varphi_a(U')=\varphi_a(U)+\langle\varphi_a(x)\rangle$ has dimension
$\dim\varphi_a(U)+1-\varepsilon_a$, and $\dim U'=d+1$. (b) The first statement is the
definition of admissibility. So $\dim P_a(U')$ stays $2$ on $\mathcal F(U)$ and is at most
$1+1$ elsewhere; the elements of $U'\setminus U=x+U$ have rank $6$. (c) $y+U'$ is the union of
the two cosets $y+U$ and $y+x+U$, both different from $U$. All elements of $y+U'$ have rank $6$
if and only if this holds for both cosets. For $a\in\mathcal F(U)$ the condition
$0\notin\varphi_a(y+U')$ holds if and only if it holds for both cosets, and for
$a\in\mathcal F(U')\setminus\mathcal F(U)$ it is the condition in the statement; by (b) this
covers $\mathcal F(U')$. (d) $\varphi_r(x+u)=\varphi_r(x)+\varphi_r(u)$. (e) Each nonzero
element of $U$ has exactly three nonzero points in its radical, so
$\sum_{a\neq0}(2^{\dim P_a(U)}-1)=3(2^d-1)<3\cdot255$; if all $\dim P_a(U)$ were $2$, the sum
would be $3\cdot255$.
\end{proof}


\subsection{The procedure and its completeness}
$\mathrm{Search}(U,H)$ takes a node $U$ of dimension $d$ and a group
$H\le GL(V)$ with $U^h=U$ for all $h\in H$.
\begin{enumerate}
\item If $d=8$: report $U$ if $\dim P_a(U)=2$ for all $a\neq0$, and return.
\item If $|\mathcal C(U)|<2^{8-d}-1$, return. (P2)
\item If $d\ge3$ and some $a$ with $\dim P_a(U)=1$ (resp.\ $0$) has fewer than one (resp.\ three)
cosets in $\mathrm{Cand}(U,a)$, return. (P3)
\item Choose $r$ with $\dim P_r(U)\le1$ (it exists by Lemma~\ref{lem:B-step}(e)). For $d\le2$, $r$
is the smallest point, as an integer $\sum_ia_i2^i$, with $\dim P_r(U)=1$, or with
$\dim P_r(U)=0$ if there is none; for $d\ge3$, $r$ minimises
$4\,|\mathrm{Cand}(U,a)|+[\dim P_a(U)=0]$ over the points with $\dim P_a(U)\le1$, ties going to
the smallest point. Let $H_r$ be a subgroup of $\{h\in H:hr=r\}$.
\item Let $c_1+U,\dots,c_m+U$ be representatives of the $H_r$-orbits on $\mathrm{Cand}(U,r)$
(Lemma~\ref{lem:B-action}(d)). For each $i$ let $H_i$ be a subgroup of the stabiliser of
$c_i+U$ in $H_r$, and call $\mathrm{Search}(U+\langle c_i\rangle,H_i)$.
\end{enumerate}
The search is $\mathrm{Search}(\langle J\rangle,H_1)$, where $H_1$ is the group generated by the
elements listed in Section~\ref{sec:n8search}, each of which fixes $J$. By
Lemma~\ref{lem:B-step}(b) every call is made with a node, and $H_i$ stabilises
$U+\langle c_i\rangle$ because it stabilises $U$ and $c_i+U$. The choice of $r$ and of the
subgroups $H_r,H_i$ does not affect correctness; only the orbits and the tree size depend on
them.

\begin{theorem}\label{thm:B-complete}
If a solution $W$ exists, then the search reports a node $U$ with $U=W^g$ for some
$g\in GL(V)$.
\end{theorem}
\begin{proof}
We show by induction on $d$ that there is a call $\mathrm{Search}(U_d,H_d)$ with
$\dim U_d=d$ and an element $g_d$ with $U_d\subseteq W^{g_d}$. For $d=1$, $W$ contains a form
of rank $6$, which is $J^{g}$ for some $g$ (all forms of rank $6$ are equivalent); then
$J\in W^{g^{-1}}$. Let $d<8$, $U=U_d$ and $W'=W^{g_d}$, a solution by
Lemma~\ref{lem:B-action}(c). By Lemma~\ref{lem:closure}, steps 2 and 3 do not return, and it
gives $c\in W'$ with $c+U\in\mathrm{Cand}(U,r)$. Let $c_i+U$ be the representative of its
$H_r$-orbit and $h\in H_r$ with $(c+U)^h=c_i+U$. As $U^h=U$,
$U+\langle c_i\rangle=(U+\langle c\rangle)^h\subseteq W'^h=W^{g_dh}$, so $U_{d+1}=U+\langle c_i\rangle$
and $g_{d+1}=g_dh$ satisfy the claim. For $d=8$, $U_8\subseteq W^{g_8}$ and both have dimension
$8$, so $U_8=W^{g_8}$, which satisfies (ii) and is reported in step 1.
\end{proof}
Since the search reports nothing, Theorem~\ref{thm:n8} follows. Of the programs, the proof
needs only that they carry out steps 1--5 as stated: that $\mathcal C(U)$, the dimensions
$\dim P_a(U)$ and $\mathrm{Cand}(U,r)$ are computed correctly, that every group element used
stabilises $U$ and fixes $r$, and that the representatives meet every orbit.


\subsection{How each step is computed}\label{app:compute}
\begin{enumerate}
\item \emph{Start.} $\mathcal C(U_1)$, $U_1=\langle J\rangle$, is obtained by testing all $2^{27}$
cosets of $\langle J\rangle$ (both elements must have rank $6$; there are no full points yet). It
has $41{,}891{,}328$ elements.
\item \emph{Children.} When $x+U$ is added, Lemma~\ref{lem:B-step}(a) updates $\dim P_a$ for
all $255$ points: for each $a$ one tests whether $\varphi_a(x)\in\varphi_a(U)$, using an echelon
basis of $\varphi_a(U)\subseteq V^*$ that is kept for every point. Then $\mathcal C(U')$ is
computed by Lemma~\ref{lem:B-step}(c): for each representative $y$ of a coset in $\mathcal C(U)$,
the representative of $y+x+U$ is looked up in $\mathcal C(U)$; if it is there, $y+U'$ is
admissible unless $\varphi_a(y)\in\varphi_a(U')$ for a new full point $a$.
\item \emph{Steps 2--4.} $|\mathcal C(U)|$ is the size of this set, and for each point $a$ with
$\dim P_a(U)\le1$ the number $|\mathrm{Cand}(U,a)|$ is the number of representatives $y$ with
$\varphi_a(y)\in\varphi_a(U)$ (Lemma~\ref{lem:B-step}(d)).
\item \emph{Step 5.} The orbits of $H_r$ on $\mathrm{Cand}(U,r)$ are found by breadth-first
search: the generators of $H_r$ act linearly on $\F_2^{28}$ ($B^g=g^{\mathsf T}Bg$), and for each
candidate the search records a product $u$ of generators mapping the orbit representative onto
it; the representative of an orbit is its first element in the order in which the candidates
are stored. Generators of $H_r$ and of $H_i$ are obtained by composing a random product $g$ of
generators of $H$ with the inverse of the recorded element $u$ that has the same image as $g$,
and as Schreier generators. Every generator is checked to stabilise $U$ and to fix $r$ (for $H_r$) or
the coset $c_i+U$ (for $H_i$); then so does every element of the group they generate.
\item \emph{Order of work.} The tree is traversed depth first. The random choices in step 5 are
determined by the path from the root, so the $486$ subtrees below the nodes of dimension $3$
can be processed separately and distributed over several processes. The numbers of nodes of
the first program at dimensions $4$ and $5$ depend on the subgroups found in step 5, hence on
these choices; those of the second program depend only on the $486$ nodes of dimension $3$.
\end{enumerate}
In the first program (the file \texttt{crsub\_enum.c} in the directory \texttt{search\_n8/}), the
function \texttt{points\_child} implements Lemma~\ref{lem:B-step}(a),(b), the function
\texttt{make\_child} implements (c), the function \texttt{choose\_point} implements steps 1--4,
and the function \texttt{expand} implements step 5. The second program carries out the same steps in Python with \texttt{numpy}; it was
written separately, from a specification of steps 1--5, and shares no code with the first.

\subsection{Symmetry and the certificate}
The symmetry reduction matters most at the top: the $740{,}736$ candidates of $U_1$ fall into
$3$ orbits, and the $727{,}936$ candidates of the three nodes of dimension $2$ into $486$ orbits.
For these four nodes the first program writes a certificate: for every candidate $x+U$, the
index $o$ of its orbit representative $c_o$ and a matrix $u\in GL(8,2)$ with $U^u=U$, $ur=r$ and
$(c_o+U)^u=x+U$. Checking it requires only linear algebra: check that $r$ is the point prescribed by step 4 (in
particular $\dim P_r(U)\le1$); compute $\mathrm{Cand}(U,r)$ from the definitions and compare
it with the listed cosets; for each line, check that $u$ is invertible, maps a basis of $U$
into $U$, fixes $r$ and maps $c_o$ into $x+U$; and check that every $c_o+U$ is a candidate and
that the nodes of the next dimension are the subspaces $U+\langle c_o\rangle$. Then, as in the proof of
Theorem~\ref{thm:B-complete}, every solution containing $U$ is equivalent to one containing
some $U+\langle c_o\rangle$. Below dimension $3$ the second program uses no symmetry at all.

\subsection{What the proof depends on}\label{app:depend}
\begin{center}
\begin{tabular}{>{\raggedright\arraybackslash}p{0.36\textwidth}>{\raggedright\arraybackslash}p{0.56\textwidth}}\toprule
Ingredient & How it is established\\\midrule
Reduction to (i) and (ii) & Theorem~\ref{thm:constrank} and Section~\ref{sec:n8search} (proved)\\
Lemmas~\ref{lem:closure}, \ref{lem:B-action}, \ref{lem:B-step} and Theorem~\ref{thm:B-complete} & proved\\
Table of forms of rank $6$ & computed twice (Gaussian elimination, Pfaffians), identical; the
count agrees with the closed formula\\
Reduction to the $486$ nodes of dimension~$3$ & certificate of $1{,}468{,}672$ matrices, checked
by the second program\\
No node of dimension $6$ below them & first program (with symmetry; two complete runs
from the archived sources) and second program (without symmetry; one complete run)\\\bottomrule
\end{tabular}
\end{center}
To verify Theorem~\ref{thm:n8} without the first program, one needs the proved statements, one
of the two rank tables, the certificate check and the search of the second program. The last
item cannot be replaced by a short certificate: it visits more than $10^9$ nodes.

\subsection{A small example}
For $n=6$ the same procedure (rank $4$, dimension $6$, $J=e_0\wedge e_1+e_2\wedge e_3$ with
radical $\langle e_4,e_5\rangle$) gives a tree that can be followed by hand. $\mathcal C(U_1)$ has
$5{,}086$ elements. Step 4 chooses $r=e_4$, the smallest point with $\dim P_r(U_1)=1$; $366$
cosets are candidates and fall into $2$ orbits, giving $2$ nodes of dimension $2$. These have
$144$ candidates in $6$ orbits. Of the $6$ nodes of dimension $3$, two fail (P3) and the others
have $32$ candidates in $15$ orbits. Of the $15$ nodes of dimension $4$, six fail (P2) and the
others have $27$ candidates, each forming its own orbit. Each of the $27$ nodes of dimension $5$
has exactly one candidate, and the $27$ nodes of dimension $6$ all satisfy (ii). They are $9$
distinct subspaces, in the two $GL(6,2)$-classes of Section~\ref{sec:n8search}. Without symmetry
reduction the procedure gives $12{,}288$ subspaces containing $J$, the same set as the plain
enumeration of Section~\ref{sec:n8search}, which does not prescribe radical points.

\subsection{Why the search is small}\label{app:small}
Let $p=149{,}920{,}960/2^{28}\approx0.5585$ be the proportion of forms of rank $6$. If the ranks of
the $2^d$ elements of a coset of a $d$-dimensional subspace behaved independently, there would be
about $2^{28-d}p^{2^d}$ cosets consisting of forms of rank $6$: about $4.2\cdot10^7$,
$6.5\cdot10^6$, $3.2\cdot10^5$, $1.5\cdot10^3$ and $0.07$ for $d=1,\dots,5$. For four random
subspaces of constant rank $6$ of each dimension the exact numbers are within $2\%$ of this for
$d\le3$, about twice as large for $d=4$ (about $3{,}200$), and between $6$ and $8$ for $d=5$; no
random $6$-dimensional subspace had any such coset. A candidate must moreover contain an element with
the chosen point $r$ in its radical: $\varphi_r(x)$ must lie in the subspace $\varphi_r(U)$, of
dimension $d-\dim P_r(U)$, of the $7$-dimensional space of linear forms vanishing at $r$, a
condition of probability about $2^{d-\dim P_r(U)-7}$; and admissibility removes further cosets
once full points appear. In the search of the second program, which uses no symmetry there, the nodes of dimension $3$
have about $17{,}700$
candidates on average and the nodes of dimension $4$ about $160$, and every node of dimension
$5$ has fewer than $7$ admissible cosets or fails (P3). This explains why the search ends at
dimension $5$; it is an explanation, not part of the proof.

\subsection{Counts}
Numbers of nodes by dimension; in parentheses the nodes returned in step 2
and in step 3. The second program takes the nodes of dimension at most $3$ from the certificate
and uses no symmetry reduction below them. In both programs no node of dimension at most $4$ is
returned in step 2 or 3.
\begin{center}
\begin{tabular}{crr}\toprule
dimension & first program & second program\\\midrule
$1$ & $1$ & $1$ (certificate)\\
$2$ & $3$ & $3$ (certificate)\\
$3$ & $486$ & $486$\\
$4$ & $1{,}466{,}318$ & $8{,}604{,}140$\\
$5$ & $244{,}686{,}904$ ($244{,}668{,}896$; $18{,}008$) & $1{,}381{,}518{,}812$ ($1{,}381{,}416{,}881$; $101{,}931$)\\
$6$ & $0$ & $0$\\\bottomrule
\end{tabular}
\end{center}

\begin{thebibliography}{99}
\bibitem{ACPBSN25} S.~Andreoli, E.~Piccione, L.~Budaghyan, P.~St\u{a}nic\u{a}, S.~Nikova.
On decompositions of permutations in quadratic functions. J.\ Cryptology 38, 22 (2025).
\bibitem{BNNRS15} B.~Bilgin, S.~Nikova, V.~Nikov, V.~Rijmen, N.~Tokareva, V.~Vitkup. Threshold
implementations of small S-boxes. Cryptogr.\ Commun.\ 7(1), 3--33 (2015).
\bibitem{BL10} C.~Bracken, G.~Leander. A highly nonlinear differentially 4 uniform power
mapping that permutes fields of even degree. Finite Fields Appl.\ 16(4), 231--242 (2010).
\bibitem{BTT12} C.~Bracken, C.H.~Tan, Y.~Tan. Binomial differentially 4 uniform permutations with
high nonlinearity. Finite Fields Appl.\ 18(3), 537--546 (2012).
\bibitem{BB64} R.H.~Bruck, R.C.~Bose. The construction of translation planes from projective
spaces. J.\ Algebra 1, 85--102 (1964).
\bibitem{Carlet21} C.~Carlet. Boolean Functions for Cryptography and Coding Theory. Cambridge
University Press (2021).
\bibitem{CP19} P.~Charpin, J.~Peng. New links between nonlinearity and differential uniformity.
Finite Fields Appl.\ 56, 188--208 (2019). doi:10.1016/j.ffa.2018.12.001.
\bibitem{DMB19} L.~De Meyer, B.~Bilgin. Classification of balanced quadratic functions. IACR
Trans.\ Symmetric Cryptol.\ 2019(2), 169--192 (2019). doi:10.13154/tosc.v2019.i2.169-192.
\bibitem{DG75} P.~Delsarte, J.-M.~Goethals. Alternating bilinear forms over GF(q). J.\ Combin.\
Theory Ser.\ A 19, 26--50 (1975).
\bibitem{Gold68} R.~Gold. Maximal recursive sequences with 3-valued recursive cross-correlation
functions. IEEE Trans.\ Inf.\ Theory 14, 154--156 (1968).
\bibitem{Gow18} R.~Gow. Connections between rank and dimension for subspaces of bilinear forms.
arXiv:1801.07529 (2018).
\bibitem{GQ09} R.~Gow, R.~Quinlan. Galois extensions and subspaces of alternating bilinear forms
with special rank properties. Linear Algebra Appl.\ 430(8--9), 2212--2224 (2009).
\bibitem{GM23} A.~Gupta, S.~Mandal. Constant rank subspaces of alternating bilinear forms from
Galois theory. arXiv:2310.03340 (2023).
\bibitem{Hir85} J.W.P.~Hirschfeld. Finite Projective Spaces of Three Dimensions. Oxford University
Press (1985).
\bibitem{Kissat} A.~Biere et al. Kissat SAT solver. \url{https://github.com/arminbiere/kissat}.
\bibitem{KMCLLJ22} K.H.~Kim, S.~Mesnager, J.H.~Choe, D.N.~Lee, S.~Lee, M.C.~Jo. On permutation
quadrinomials with boomerang uniformity 4 and the best-known nonlinearity. Des.\ Codes Cryptogr.\
90, 1437--1461 (2022).
\bibitem{KP} L.~K\"olsch, A.~Polujan. The combinatorial structure and value distributions of
plateaued functions. J.\ Cryptology 39, 21 (2026).
\bibitem{MTX19} S.~Mesnager, C.~Tang, M.~Xiong. On the boomerang uniformity of quadratic
permutations. Des.\ Codes Cryptogr.\ 88, 2233--2246 (2020).
\bibitem{NNR19} S.~Nikova, V.~Nikov, V.~Rijmen. Decomposition of permutations in a finite field.
Cryptogr.\ Commun.\ 11, 379--384 (2019).
\bibitem{NRR06} S.~Nikova, C.~Rechberger, V.~Rijmen. Threshold implementations against side-channel
attacks and glitches. ICICS 2006, LNCS 4307 (2006). doi:10.1007/11935308\_38.
\bibitem{Nyb93} K.~Nyberg. Differentially uniform mappings for cryptography. EUROCRYPT '93, LNCS 765,
55--64 (1994).
\bibitem{QXL13} L.~Qu, H.~Xiong, C.~Li. A negative answer to Bracken--Tan--Tan's problem on
differentially 4-uniform permutations over $\F_{2^n}$. Finite Fields Appl.\ 24, 55--65 (2013).
\bibitem{WHH14} N.~Wetzler, M.J.H.~Heule, W.A.~Hunt Jr. DRAT-trim: efficient checking and trimming
using expressive clausal proofs. SAT 2014, LNCS 8561 (2014). doi:10.1007/978-3-319-09284-3\_31.
\end{thebibliography}
\end{document}
